% -*- program: xelatex -*-
 
% Dr Driver's standard reading quiz
% Source: https://gist.github.com/danieldriver/90a73c4d3c72dd837e39#file-quiz-tex 
\documentclass[%
% 11pt,
% answers,
addpoints]{exam}
 
\pagestyle{head}
\firstpageheader{% left
	BSTH \liningnums{314}: Psalms --- Reading Quiz 5 \\
	Dr. Daniel R. Driver, 19 August 2014
}{% center - blank
}{% right
	\ifprintanswers \textbf{Answer Key}
	\else Full name:\enspace\makebox[2in]{\hrulefill} \fi \\
	This quiz is worth a total of \textbf{\numpoints\ points}.
}
\runningheader{}{}{}
% For double-sided quizzes
% \pagestyle{headandfoot}
% \footer{}{\thepage}{}
 
% Typography and layout
\usepackage{fontspec,realscripts}
\defaultfontfeatures{Ligatures=TeX}
%\setmainfont{Meta Serif Pro}
%\setsansfont{Meta Pro}
\frenchspacing
% - print solutions in sans serif
\unframedsolutions
\SolutionEmphasis{\sffamily}
\renewcommand{\solutiontitle}{}
% - box points & center in the right margin w/ custom setup@point@toks
\boxedpoints
\pointsinrightmargin
\marginbonuspointname{\textsc{up}}
\makeatletter% rewrite setup@point@toks assuming right margins
\def\clap#1{\hbox to 0pt{\hss#1\hss}}% define \clap as per https://www.tug.org/TUGboat/tb22-4/tb72perlS.pdf
\def\setup@point@toks{%
	\point@toks={%
		\rlap{\hskip-\@totalleftmargin
			  \hskip\textwidth
			  \hskip\@rightmargin
			  \hskip-\rightpointsmargin
			  \clap{\padded@point@block}% change \llap to \clap
		}%
		\global \point@toks={}%
	}%
}% end setup@point@toks
\setlength{\rightpointsmargin}{.5in}% assuming the default 1" margins
\makeatother
% - adjust the top and bottom margins
\extraheadheight{.25in}
\extrafootheight{-.5in}
\setlength{\marginparwidth}{1.5in}
% NB: remember to use \newpage after the last question
 
\begin{document}
 
	% Quiz instructions
	\begin{center}
	\parbox{5.5in}{\centering\itshape Answer each of the following questions in one or two complete sentences.}
	\end{center}
 
	% Quiz questions
	\begin{questions}
 
	\question[2]
		Question
 
		\begin{solution}[\stretch{1}]
			Answer
		\end{solution}
 
	\question[2]
		Question
 
		\begin{solution}[\stretch{1}]
			Answer
		\end{solution}
 
	\question[2]
		Question
 
		\begin{solution}[\stretch{1}]
			Answer
		\end{solution}
 
	\question[2]
		Question
 
		\begin{solution}[\stretch{1}]
			Answer
		\end{solution}
 
	\question[2]
		Question
 
		\begin{solution}[\stretch{1}]
			Answer
		\end{solution}
 
	\bonusquestion[2]
		Question
 
		\begin{solution}[\stretch{1}]
			Answer
		\end{solution}
 
	\newpage
	\end{questions}
\end{document}